A alocação de slots das companhias aéreas dentro da infraestrutura aeroportuária influencia diretamente a lucratividade de um aeroporto e sua eficiência em absorver o crescente tráfego de transporte aéreo de passageiros. Embora a expansão da capacidade aeroportuária represente a solução definitiva para o desequilíbrio entre esse crescimento e a infraestrutura limitada, fatores como regulamentações, áreas delimitadas e limitações financeiras frequentemente exigem que os aeroportos implementem mecanismos alternativos de gestão da demanda.
%
Essa dissertação tem como objetivo apresentar um método para resolver o problema de alocação de slots aeroportuários, focando em um modelo que emprega Programação Inteira (\textit{Integer Programming} -- \acrshort{IP}), Aprendizado de Máquinas (\textit{Machine Learning} -- \acrshort{ML}) e Teoria das Filas (\textit{Queueing Theory} -- \acrshort{QT}). 
%
<<UNIFICAR ESSAS TRES FRASES SEGUINTES(SOBRE AS TÉCNICAS QT, ML E IP) EM UMA ÚNICA FRASE>> Especificamente, \acrshort{QT} é utilizada para simular a capacidade real de um aeroporto, identificando gargalos. \acrshort{ML}, por sua vez, busca preencher lacunas de informação cruciais para um processo de alocação aprimorado, particularmente detalhes frequentemente retidos pelas companhias aéreas para projeções de longo prazo. \acrshort{IP} tem sido amplamente adotada na literatura como um mecanismo de gestão de slots, levando à redução de atrasos e deslocamento de slots para as companhias aéreas. 
%
%Adicionalmente, uma Literatura Sistemática de Literatura (\acrshort{SLR}) foi conduzida, investigando as abordagens existentes para o problema de alocação de slots. Este processo gerou um protocolo de revisão, dados de pesquisa e um relatório abrangente detalhando o conhecimento avaliado sob múltiplas perspectivas, as metodologias empregadas e as áreas identificadas para pesquisa futura.
%
Essa dissertação busca oferecer uma contribuição tríplice ao propor um método que integre as técnicas de QT, ML e IP para a alocação eficiente de \textit{slots} em aeroportos, apresentar um estudo de caso aplicado ao Aeroporto Internacional de Brasília (\acrshort{BSB}) para avaliar sua capacidade de operação pelo emprego de \acrshort{QT} e, por fim, propor uma alocação de submissões iniciais de \textit{slots} solicitados pelas companhias aéreas pela utilização de \acrshort{IP} baseado em informações faltantes obtidas por meio da aplicação de \acrshort{ML}.
%
